
On
trouve donc un homomorphisme de spécialisation sur @@M0@@, mais à valeurs non dans @@M1@@, mais dans @@M2@@, où @@M3@@ est déduit de @@M4@@ par une extension @@M5@@ des corps de base (qu'on peut supposer extension de type fini si @@M6@@ est universellement japonais). Passant aux fibres géométriques comme expliqué dans 7.16, on trouve de même un homomorphisme
@@M7@@
donc ici encore, l'homomorphisme de spécialisation «~géométrique~» prend ses valeurs non dans @@M8@@, mais dans @@M9@@, où @@M10@@ est déduit de @@M11@@ par l'extension de corps algébriquement clos @@M12@@.

\textbf{7.17.3.} En général, l'homomorphisme @@M13@@ n'est pas un isomorphisme (à cause de la présence possible de composantes «~continues~» dans le groupe @@M14@@). Cependant, il est facile de vérifier que c'est un isomorphisme modulo équivalence algébrique (7.11), et aussi modulo équivalence numérique (4.5.1) si @@M15@@ est lisse et propre. Par suite, utilisant 7.12, on trouve que (7.17.2) induit un homomorphisme de spécialisation
@@M16@@
et si @@M17@@ est propre et lisse, et sous réserve que l'homomorphisme de spécialisation (7.17.2) soit compatible avec l'équivalence numérique (cf. 7.9.4), on trouve de même un homomorphisme
@@M18@@
Nous laissons au lecteur le soin de développer (en cas de besoin) les constructions correspondantes pour @@M19@@, @@M20@@, et @@M21@@, et les propriétés de ces homomorphismes de spécialisation qui résultent trivialement des propriétés développées dans les sections précédentes. Signalons seulement
